\markdownRendererDocumentBegin
\markdownRendererWarning{The `hybrid` option has been soft-deprecated.}{The `hybrid` option has been soft-deprecated.}{Consider using one of the following better options for mixing TeX and markdown: `contentBlocks`, `rawAttribute`, `texComments`, `texMathDollars`, `texMathSingleBackslash`, and `texMathDoubleBackslash`. For more information, see the user manual at <https://witiko.github.io/markdown/>.}{Consider using one of the following better options for mixing TeX and markdown: `contentBlocks`, `rawAttribute`, `texComments`, `texMathDollars`, `texMathSingleBackslash`, and `texMathDoubleBackslash`. For more information, see the user manual at <https://witiko.github.io/markdown/>.}\markdownRendererSectionBegin
\markdownRendererHeadingOne{Rootear un árbol}\markdownRendererInterblockSeparator
{}Rootear un árbol consiste en designar un nodo como raíz y establecer relaciones padre-hijo mediante DFS o BFS desde ese nodo. Este proceso transforma un árbol no dirigido en un árbol dirigido con jerarquía, facilitando operaciones como cálculo de ancestros, profundidades y aplicación de algoritmos de programación dinámica en árboles. La complejidad es $O(n)$ y la elección de la raíz puede optimizarse usando el centroide para minimizar la altura.\markdownRendererInterblockSeparator
{}\markdownRendererInputFencedCode{./_markdown_main/88da1d3c1e59827508471827f28b8f96.verbatim}{cpp}{cpp}
\markdownRendererSectionEnd \markdownRendererDocumentEnd